\section{Finite Fano data and their geometric scope}
\label{sec:fano-data}
The numerical classification uses nine small matrices and eight classical
rationality assertions. The matrix calculation establishes the spectral
conditions; the cited family identifications and smooth deformation bridges
establish their every-member scope.

\subsection{Matrices and residues}
In the classical anticanonical-power basis, up to a scalar Euler shift,
write the matrix at Novikov parameter one as
\[
 U=\begin{pmatrix}a&c&e&f\\1&b&d&e\\0&1&b&c\\0&0&1&a\end{pmatrix},
 \qquad D=\tfrac12\diag(3,1,-1,-3).
\]
The nine inputs are:
\begin{center}\small
\begin{tabular}{@{}lrrrrrr@{}}
\toprule
Family & $a$ & $b$ & $c$ & $d$ & $e$ & $f$\\\midrule
g=2 & 120&744&137520&650016&119681280&21690374400\\
g=3 & 24&104&3888&13600&504576&18323712\\
g=4 & 12&42&792&2340&43632&793152\\
g=5 & 8&24&304&800&9984&121088\\
g=6 & 6&16&156&380&3600&33120\\
g=8 & 4&9&64&140&924&5936\\
d=1 & 0&0&240&1248&0&57600\\
d=2 & 0&0&48&160&0&2304\\
d=3 & 0&0&24&60&0&576\\
\bottomrule
\end{tabular}
\end{center}
The columns $(1,U1,U^21,U^31)$ have determinant one, so the matrix is cyclic.
For genera $2,3,4,5$ the characteristic polynomial is $T^3(T-s)$, with
$s=1728,256,108,64$, respectively. For genera $6,8$ it is
$T^2(T^2-44T-16)$ and $T^2(T-27)(T+1)$.
For degrees $1,2,3$ it is $T^2(T^2-s)$, with $s=1728,256,108$.
Thus the displayed repeated factors are whole primary factors.
Restoring the anticanonical-degree parameter $t$ gives
$tB(t)UB(t)^{-1}$, where $B(t)=\diag(1,t^{-1},t^{-2},t^{-3})$;
in index two the integral Novikov variable is $q=t^2$.
This homogeneous substitution, after inverting $t$, preserves primary ranks
and the canonical residue discriminant. A scalar Euler shift changes only
the removed exponential.

For a repeated rank-two factor, choose a Jordan chain and a complementary
invariant frame. In that frame let $J$ be the leading matrix and $B$ the
grading coefficient. Solve the off-diagonal Sylvester equation
$[J,X]+B_{\rm off}=0$ for the first regular gauge coefficient.
The next diagonal coefficient is $(BX)_{ii}$: the derivative term $-X$
is off-diagonal. Applying the elementary modification to this finite jet
gives exactly the five discriminants in Proposition~\ref{prop:fano-endpoints}.
The independent checker constructs the frame by rational elimination;
the other checker computes the projector and reduced inverse directly.
Neither requires a higher-rank logarithmic lattice.

\subsection{Sources and all smooth members}
The counting matrices agree entrywise with the reconstructions accompanying
B\"ohning--Graf von Bothmer--Su'a \cite[Sections~9--10 and associated data]{NaiveAtoms}.
If $M$ denotes a source matrix at parameter one, the displayed matrix is
$D_0^{-1}M^{\mathsf t}D_0+aI$. For index one of degree $k=2g-2$,
$D_0=\diag(1,1,k,k)$; for index two of degree $k$,
$D_0=\diag(1,2,4k,8k)$ and $a=0$.
The input period coefficients through degree eight agree with the formulas
and tables in Coates--Corti--Galkin--Kasprzyk
\cite[Sections~3--17]{CCGK}. The genus-six and genus-eight quantum
inputs are proved in \cite[Theorem~6.1.1]{PrzyjalkowskiGenera}; the weighted
and complete-intersection inputs are identified in the cited period tables.
The reconstruction constrains the quantum-product matrix by degree, classical
products and pairing, then determines its coefficients from the period equation.
The finite coefficient comparison checks normalization; the GW identification
is the imported reconstruction result.

The smooth classification has seventeen deformation families, indexed by
index and degree or genus; we use its statement and Hodge numbers as recalled
in \cite[Introduction and Table~5]{KuznetsovProkhorovNodal}, which attributes
them to Iskovskikh--Prokhorov. The family presentations in
\cite[Sections~3--17]{CCGK} include the weighted models.
Smooth deformation invariance transports the quantum products with the ample
generator marked. It does not transport rationality or identify Hodge
structures of different members.

Two presentations require an explicit bridge. The special genus-three model
and the quartic are smooth intersections of type $(2,4)$ in
$\PP(1,1,1,1,1,2)$ \cite[Section~9]{CCGK}. Perturbing the degree-two equation
to have nonzero coefficient of the weight-two coordinate gives a quartic
by elimination. Smoothness is open and the smooth parameter locus is
connected. For a special Gushel--Mukai threefold, its linear section of the
cone passes through the vertex, while its smooth quadric section avoids the
vertex \cite[Sections~2--3]{DebarreGM}. Perturb the linear section to miss
the vertex. Properness and openness of smoothness give a smooth local family
with ordinary general member; ordinary presentations form an open subset of
an irreducible parameter space. Thus both special strata have the same
small quantum matrix as their ordinary smooth family.

The remaining eight families are rational for every smooth complex member
\cite[Section~1.1]{KuznetsovProkhorov}. No quantum matrix for those controls
is a premise of the classification theorem.

\subsection{Reproduction and formal boundary}
The bundle \texttt{verification/fano-matrices/} contains the exact rational
inputs, the symbolic projector calculation, an independent rational-elimination
checker, and the source-normalization comparison. Its README specifies
commands, pinned source hashes, finite domains and trust limits.
The two implementations agree on all seventeen characteristic polynomials
and cyclic frames and all nine rank-two residues, including four zero
controls. Only the nine detected inputs above are needed here.
The computations do not establish the cited Gromov--Witten formulas,
classification, deformation invariance or geometric transport.
The manuscript's claim map records the new geometric theorems as absent
from the formalization; existing algebraic fragments retain their stated
boundaries.
